\documentclass[10pt,twocolumn]{article}
\usepackage[margin=0.72in]{geometry}
\usepackage{amsmath,amssymb,booktabs,array}
\usepackage[T1]{fontenc}
\usepackage{lmodern}
\usepackage{xcolor}
\usepackage[hidelinks,breaklinks]{hyperref}
\usepackage{microtype}
\usepackage{balance}
\usepackage{enumitem}
\usepackage{adjustbox}
\setlength{\columnsep}{0.25in}
\setlength{\emergencystretch}{1.5em}
\setlist{nosep,leftmargin=*}
\newcommand{\src}[1]{\nolinkurl{#1}}
\newcommand{\lk}[1]{\texttt{#1}}
% median[min,max] over repeats, as printed by scripts/make_tables.py
\newcommand{\rng}[2]{{\scriptsize\,[#1,\,#2]}}
% A number or result the text wants but no FINDINGS.md entry or JSON provides.
\newcommand{\todo}[1]{{\color{red}\textbf{[TODO: #1]}}}
%% generated by scripts/make_tables.py -- do not edit
\newcommand{\BurOffHi}{1.44}
\newcommand{\BurOffLo}{1.17}
\newcommand{\BurOnHi}{4.97}
\newcommand{\BurOnLo}{4.22}
\newcommand{\CbarCEight}{2887}
\newcommand{\CbarCSixteen}{2449}
\newcommand{\CbarFcRemoteAB}{4883}
\newcommand{\CbarNinetyFive}{2812}
\newcommand{\LHNinetyFive}{3.85}
\newcommand{\LocTests}{977}
\newcommand{\LocTotal}{4189}
\newcommand{\OCEight}{1189}
\newcommand{\OCSixteen}{1145}
\newcommand{\OFcOld}{784}
\newcommand{\OFcRemoteAB}{937}
\newcommand{\ONeededEighty}{703}
\newcommand{\PromotedCEight}{0.216}
\newcommand{\SusOffHi}{1.55}
\newcommand{\SusOffLo}{1.50}
\newcommand{\SusOnHi}{1.63}
\newcommand{\SusOnLo}{1.58}
\newcommand{\UtilAtOFcOld}{0.782}
\newcommand{\UtilCEight}{0.708}
\newcommand{\UtilCSixteen}{0.681}
\newcommand{\UtilFcOld}{0.860}
\newcommand{\UtilFcRemoteAB}{0.839}
\newcommand{\XNinetyFive}{0.627}


\title{Scheduler Subversion by Delegation Locks in Cooperative Executors:\\
Combiner Burden, Service Unfairness, and Lock-Side Policies}
\author{Draft --- \texttt{crates/coro\_delegation}}
\date{}

\begin{document}
\maketitle

\begin{abstract}
Delegation (combining) locks execute other tasks' critical sections on one
thread. In a cooperatively scheduled, work-stealing coroutine executor this
subverts the scheduler in two ways. First, \emph{combiner burden}: the worker
that becomes combiner pays for everyone's critical sections, and so does every
task queued behind it. In our measurements, a Combine-and-Exchange (CES) lock
keeps one worker combining for the whole 2\,s window (burden Jain index $1/W$),
and, with executor balancing disabled, flat combining (FC) without a
cooperative yield lets one of 64 clients
perform every operation. Second, \emph{service unfairness}: FIFO hand-off gives
each client service in proportion to its critical-section cost; with a 1:8
cost mix every FIFO lock we measured, \lk{tokio::sync::Mutex} included, has the
service Jain index predicted by that proportionality (0.65--0.66 sustained). We treat
\emph{where the combiner runs} and \emph{whom it serves next} as lock policies
that need only wake-placement hints from the executor. A 64-hand-off chain
bound with a home-worker chain break (CES) and a pass-end hand-off with remote
wakes (FC) restore burden Jain 0.98--1.00 at 0.95--1.01$\times$ CES
throughput. Usage-ordered FC with pass limit 16, home wakes and a 16-pass
starvation clamp reaches service Jain 0.997 with a worst queue wait of
17 passes. Under sustained load the two burden-fair locks run
\SusOnLo--\SusOnHi$\times$ \lk{tokio::sync::Mutex} (tokio 1.53.1), and
\SusOffLo--\SusOffHi$\times$ with tokio's LIFO slot disabled. A per-lock
server task (actor) reaches either delegation throughput or burden fairness,
not both. Fairness costs lock utilisation: the lock's per-operation overhead of
about 1\,150 cycles is paid by cheap operations too, so utilisation falls from
0.84 to 0.68.
\end{abstract}

% ---------------------------------------------------------------------------
\section{Introduction}\label{sec:intro}

Coroutine runtimes such as tokio~\cite{tokio} multiplex many tasks onto a few
worker threads. Scheduling is cooperative: a task runs until it returns
\lk{Pending}. The runtime balances load by work stealing. Its fairness
guarantees are weak and are stated only for runnable tasks. Locks for such
runtimes are \emph{async mutexes}: a waiter suspends, and the unlocker wakes
it through the scheduler. CES~\cite{ces} observed that this hand-off puts a
scheduler round trip on the critical path. CES instead suspends the unlocker
and resumes the next waiter inline on the unlocking thread, so contended
critical sections stay on one thread. This is a delegation lock in the sense
of flat combining (FC)~\cite{fc} and TCLocks~\cite{tclocks}. CES reports
throughput only.

Patel et al.~\cite{scl} showed that ordinary locks \emph{subvert} the OS
scheduler: lock hold time, not the scheduler's share, decides who gets CPU
time. We find two analogous subversions when delegation locks meet a
cooperative executor:

\begin{enumerate}
\item \textbf{Combiner burden.} Critical sections run on whichever worker is
  combining, and the executor cannot see that. With CES the unlocking worker
  stays combiner for as long as the wait queue is non-empty. Under sustained
  contention that is the entire 2\,s run: the burden Jain index over workers
  is exactly $1/W$ in every saturated sustained cell at $W{=}4,8,16$ (\S\ref{sec:motiv}), and a bystander task
  queued on that worker is never polled. Without executor balancing, FC without a cooperative
  yield is worse. The combining task never returns \lk{Pending}, so 63 of 64 clients do
  no work for 2\,s.
\item \textbf{Service unfairness.} FIFO hand-off gives equal operations per
  client, so lock time is proportional to critical-section cost. For a 1:8
  light:heavy cost mix, the resulting service Jain index is 0.653 (from the
  measured 1\,306:8\,324 cycles per op). \lk{dispatch}, CES, FC and
  \lk{tokio::sync::Mutex} all reproduce it to within 0.01 (Table~\ref{tab:fifo}).
\end{enumerate}

Our position is that both are lock policies, not executor policies. A
delegation lock decides \emph{where the combiner runs} and \emph{whom it serves
next}. To act on the first decision, the only thing it needs from the executor
is a placement argument on wakes (inline, remote, home). We make three
contributions:

\begin{itemize}
\item Measurements of both subversions on a work-stealing executor built for
  this study (\S\ref{sec:motiv}). They cover \lk{dispatch}, CES, FC, a
  usage-ordered FC (FC-PQ), and a cross-check against the real tokio locks.
\item A minimal executor contract (\lk{wake\_with(placement)},
  \lk{reschedule\_self\_remote}, a run-next slot), and lock policies on top of
  it (\S\ref{sec:design}). For burden: a CES chain bound with home-worker chain
  break, and FC yield-after-combine with remote wakes. For service: usage-ordered
  FC-PQ with pass limit 16 and a starvation clamp.
\item An evaluation (\S\ref{sec:eval}) of what these policies buy and cost.
  \lk{ces-k64-home} and \lk{fc-remote} restore burden Jain 0.98--1.00 at
  0.95--1.01$\times$ CES throughput. \lk{fcpq-h16-home-c16} reaches service Jain
  0.997. Burden-fair delegation runs \SusOnLo--\SusOnHi$\times$
  \lk{tokio::sync::Mutex} under sustained load; its 4--5$\times$ bursty margin
  is mostly an artefact of tokio's LIFO slot. The obvious alternative, a
  per-lock server task (actor), reaches throughput or burden fairness but not
  both. Service fairness lowers lock utilisation by an amount that the
  per-op overhead $o$ fully accounts for.
\end{itemize}

We do not claim the first delegation lock for cooperative scheduling. CES
states that claim for itself~\cite{ces}. Our subject is fairness, and the
executor/lock split that it requires.

% ---------------------------------------------------------------------------
\section{Background and Motivation}\label{sec:motiv}

\subsection{Cooperative executors}
A tokio-style executor~\cite{tokio} runs $W$ worker threads. Each has a local
run queue, and there is one shared injector queue. An idle worker steals half
of a random peer's queue. tokio adds two policies that matter here:
\begin{itemize}
\item A \emph{LIFO slot}: a task woken from a worker is polled right after the
  current poll, cannot be stolen, and is capped at 3 in a row.
\item A \emph{cooperative budget} of 128 units per poll, consumed by
  \lk{tokio::sync::Mutex} acquires.
\end{itemize}
Tokio's fairness statement for its mutex covers only the tasks waiting for
it~\cite{tokio6049}.

Our executor (\S\ref{sec:impl}) is built from \lk{async-task}~\cite{asynctask}
and \lk{crossbeam-deque}~\cite{crossbeam}. It has per-worker FIFO queues, an
injector and a single run-next slot. Every 31 polls a worker drains up to
$\mathit{len}/W{+}1$ injector tasks (tokio's share rule). In our workload every
worker hosts an always-runnable bystander, so steal-when-idle never fires. We
therefore add a rate-limited \emph{balancing steal}: every $b$ polls, a busy
worker takes half of a random peer's queue. Runs use $b{=}31$ (``b31'') or no
balancing (``b0'').

\subsection{Locks}
We use the taxonomy of CES~\cite{ces}.
\begin{itemize}
\item \emph{Dispatch} (\lk{dispatch}; tokio, Kotlin, Boost): unlock enqueues
  the next waiter and the owner continues.
\item \emph{Inline}: unlock resumes the waiter recursively.
\item \emph{CES} (\lk{ces}): unlock reschedules the owner remotely and resumes
  the head waiter inline.
\item \emph{FC} (\lk{fc})~\cite{fc}: clients publish closures, and whoever wins
  a \lk{try\_lock} runs a pass of up to $H$ of them.
\item \emph{FC-PQ} (\lk{fcpq}): FC that serves the lowest cumulative charged
  usage first.
\end{itemize}

\subsection{Workload and metrics}
The protected structure is a \lk{BTreeMap<u64,u64>} with 65\,536 keys. A
critical section (CS) is one insert plus a TSC-timed spin: 1\,000 cycles for
light clients and $8\times$ for heavy clients, half the clients each.
\begin{itemize}
\item \emph{Sustained}: 64 clients with 4\,000 cycles of parallel work between
  operations.
\item \emph{Bursty}: 16 clients with 32\,000 cycles.
\end{itemize}
$W$ bystander tasks spin 1\,000 cycles and then yield. Three metrics are Jain
indices~\cite{jain}, $J=(\sum x_i)^2/(n\sum x_i^2)$:
\begin{itemize}
\item \emph{service Jain}, over per-client CS cycles;
\item \emph{burden Jain}, over per-worker combining cycles;
\item the \emph{combiner-worker bystander p99}, the p99 schedule-to-poll delay
  on the worker with the most combining cycles, against the maximum p99 of the
  other workers.
\end{itemize}
A bystander first polled after the window is \emph{censored} (``cens.'' in the
tables): it never ran.

\begin{table*}[t]
\centering\footnotesize
\adjustbox{max width=\linewidth}{%% generated by scripts/make_tables.py -- do not edit
\begin{tabular}{llrrrrrr}
\toprule
cell & lock & Mops/s & svc.\ $J$ & burden $J$ & comb.\ p99 & other p99 & starved c/b \\
\midrule
sus. b0 & \texttt{dispatch} & 0.209\rng{0.199}{0.210} & 0.650\rng{0.650}{0.651} & -- & -- & 5.8\rng{5.8}{6.1} & 0/0 \\
sus. b0 & \texttt{ces} & 0.404\rng{0.363}{0.413} & 0.574\rng{0.557}{0.627} & 0.125 & cens. & 2.8\rng{2.7}{2.8} & 5\rng{5}{6}/1 \\
sus. b0 & \texttt{fc-noyield} & 0.392\rng{0.391}{0.392} & 0.016 & 0.125 & cens. & 0.1 & 63/1 \\
sus. b0 & \texttt{fc} & 0.235\rng{0.234}{0.235} & 0.649 & 0.125 & 268.1 & 0.1\rng{0.1}{0.2} & 0/0 \\
sus. b0 & \texttt{fcpq} & 0.234\rng{0.231}{0.235} & 0.649\rng{0.649}{0.650} & 0.125 & 268.1 & 0.1\rng{0.1}{0.2} & 0/0 \\
\midrule
sus. b31 & \texttt{dispatch} & 0.216\rng{0.208}{0.218} & 0.661\rng{0.661}{0.664} & -- & -- & 0.8\rng{0.7}{0.8} & 0/0 \\
sus. b31 & \texttt{ces} & 0.380\rng{0.378}{0.384} & 0.651\rng{0.650}{0.652} & 0.125 & -- & 3.0\rng{3.0}{3.1} & 0/0 \\
sus. b31 & \texttt{fc-noyield} & 0.391\rng{0.390}{0.392} & 0.657\rng{0.657}{0.658} & 0.999\rng{0.998}{0.999} & 1.2\rng{1.2}{1.3} & 1.3 & 0/0 \\
sus. b31 & \texttt{fc} & 0.390\rng{0.389}{0.392} & 0.658 & 0.999 & 1.3 & 1.3 & 0/0 \\
sus. b31 & \texttt{fcpq} & 0.431\rng{0.429}{0.432} & 0.713\rng{0.713}{0.714} & 0.998\rng{0.997}{0.999} & 1.4\rng{1.3}{1.5} & 1.5\rng{1.4}{1.5} & 0/0 \\
\midrule
bur. b0 & \texttt{dispatch} & 0.057\rng{0.056}{0.057} & 0.652\rng{0.651}{0.652} & -- & -- & 18.6 & 0/0 \\
bur. b0 & \texttt{ces} & 0.315\rng{0.312}{0.318} & 0.679\rng{0.678}{0.679} & 1.000\rng{0.999}{1.000} & 44.7 & 44.7 & 0/0 \\
bur. b0 & \texttt{fc-noyield} & 0.054 & 0.062 & 0.125 & cens. & 0.1 & 15/1 \\
bur. b0 & \texttt{fc} & 0.059 & 0.650\rng{0.650}{0.651} & 0.125 & 268.1 & 0.1\rng{0.1}{0.2} & 0/0 \\
bur. b0 & \texttt{fcpq} & 0.059\rng{0.058}{0.059} & 0.651\rng{0.650}{0.651} & 0.125 & 268.1 & 0.1\rng{0.1}{0.2} & 0/0 \\
\midrule
bur. b31 & \texttt{dispatch} & 0.071\rng{0.070}{0.071} & 0.673\rng{0.672}{0.675} & -- & -- & 0.8 & 0/0 \\
bur. b31 & \texttt{ces} & 0.316\rng{0.315}{0.317} & 0.678\rng{0.678}{0.680} & 1.000\rng{0.999}{1.000} & 44.7 & 44.7 & 0/0 \\
bur. b31 & \texttt{fc-noyield} & 0.298\rng{0.296}{0.298} & 0.660\rng{0.659}{0.660} & 1.000 & 29.8\rng{28.9}{29.8} & 29.8 & 0/0 \\
bur. b31 & \texttt{fc} & 0.302\rng{0.301}{0.302} & 0.660\rng{0.660}{0.661} & 1.000 & 29.8 & 29.8 & 0/0 \\
bur. b31 & \texttt{fcpq} & 0.304\rng{0.304}{0.305} & 0.672\rng{0.671}{0.672} & 1.000 & 29.8 & 29.8 & 0/0 \\
\bottomrule
\end{tabular}
}
\caption{The problem, phase-2 matrix, $W{=}8$, heavy $8\times$ (09-28 window).
  Service $J$ is over clients and burden $J$ over workers. The p99 columns are
  bystander schedule-to-poll delays in $\mu$s on the combiner worker and the
  maximum over the other workers. The last column counts starved clients and
  bystanders (0 ops in 2\,s). \lk{fc-noyield} is FC without the cooperative
  yield of \S\ref{sec:design}.}
\label{tab:motiv}
\end{table*}

\subsection{Measured subversion}
Table~\ref{tab:motiv} shows the problem at 8 workers. The full matrix covers
$W\in\{4,8,16\}$, heavy $\in\{1,8\}$ and $b\in\{0,31\}$ (432 runs).

\paragraph{CES concentrates all combining on one worker.}
In every saturated sustained cell CES has burden Jain exactly $1/W$ (0.125 at
8 workers, 0.062 at 16), with or without balancing. The median chain is
819\,200 inline hand-offs, i.e.\ one chain lasts the whole window. What happens
to the tasks around it depends on balancing:
\begin{itemize}
\item Without balancing, CES traps whatever sits in the combiner's queue: 5
  clients (range 5--6 across repeats) and the bystander starve for 2\,s.
  Service Jain drops to 0.574\rng{0.557}{0.627}.
\item With balancing, the combiner's bystander is stolen during warm-up and
  never returns. The burden shows up as eviction, not as delay.
\end{itemize}
Whether a CES chain ever ends is an executor property. An earlier executor
version drained one injector task per 31 polls. At $W\le8$ that let the wait
queue empty periodically: burden Jain was 0.96 at 8 workers but 0.06 at 16.

\paragraph{FC without a yield hands the lock to one client.}
When the executor does not balance, the FC combiner's \lk{run} completes
synchronously. The waiters it served are woken into its own worker's queue, and
the client loop never returns \lk{Pending}. The same task therefore
re-publishes and re-wins the election forever. Service Jain is $0.016=1/64$
and 63 clients starve (bursty: 15 of 16). A yield after combining fixes the
starvation. It does not fix placement: at b0, all clients converge on the
combiner's worker (burden 1/W) and throughput is 0.235 instead of
0.390\,Mops/s with balancing.

\paragraph{FIFO service is proportional to cost.}
Table~\ref{tab:fifo} applies the two-class model
$J(x)=(1+x)^2/(2(1+x^2))$ to each lock's measured costs, where $x$ is the
light:heavy per-client service ratio. Service Jain follows from the costs to
three digits for \lk{dispatch}, CES, FC and \lk{tokio::sync::Mutex}. Equal
service would need 6.0--6.5 light operations per heavy one.

\begin{table}[t]
\centering\footnotesize\setlength{\tabcolsep}{3pt}
\adjustbox{max width=\linewidth}{%% generated by scripts/make_tables.py -- do not edit
\begin{tabular}{lrrrrrr}
\toprule
lock & CS$_L$ & CS$_H$ & L:H ops & model $J$ & measured $J$ & L:H for $J{=}1$ \\
\midrule
\texttt{dispatch} & 1391 & 8403 & 1.00 & 0.661 & 0.661\rng{0.661}{0.664} & 6.04 \\
\texttt{ces} & 1280 & 8295 & 1.00 & 0.651 & 0.651\rng{0.650}{0.652} & 6.48 \\
\texttt{fc} & 1354 & 8355 & 1.00 & 0.658 & 0.658 & 6.17 \\
\texttt{tokio-mutex} & 1404 & 8402 & 1.00 & 0.663 & 0.663\rng{0.662}{0.663} & 5.98 \\
\bottomrule
\end{tabular}
}
\caption{FIFO service fairness, $W{=}8$, sustained, b31 (09-28 window).
  CS$_{L,H}$ are critical-section cycles per op (TSC). The model uses
  $x=\text{L:H}\cdot\text{CS}_L/\text{CS}_H$. The last column is
  CS$_H$/CS$_L$.}
\label{tab:fifo}
\end{table}

\paragraph{Off-the-shelf async locks.}
In tokio, only \lk{tokio::sync::Mutex} serves every client (\S\ref{sec:q3}).
\lk{async-lock}~\cite{asynclock} collapses to a single client in 11 of 12
matrix runs (63 of 64 or 15 of 16 clients starved). With the LIFO slot on it
collapses in 20 of 24 runs, and with the slot off in 0 of 12. \lk{std} and
\lk{parking\_lot} mutexes block the worker thread. The first $W$ clients to be
polled keep the $W$ workers for the whole run, so starved clients equal clients
$-\,W$, and all $W$ bystanders starve.

% ---------------------------------------------------------------------------
\section{Design Goals}\label{sec:goals}

Following SCL~\cite{scl}, we state the properties a delegation lock in a
cooperative executor should have.
\begin{enumerate}[label=\textbf{G\arabic*}]
\item \textbf{Burden fairness.} Combining cycles are spread across workers
  (burden Jain $\ge0.9$). The p99 of a bystander on the combiner worker is
  within one histogram bucket of the p99 elsewhere.
\item \textbf{Service fairness.} Clients receive equal lock \emph{time}, not
  equal operations. This requires usage accounting (service Jain $\ge0.95$
  under 1:8 cost heterogeneity).
\item \textbf{Work conservation.} Throughput is no worse than FIFO
  delegation. The fairness policy adds no idle lock time.
\item \textbf{Executor independence.} The lock asks the executor only for
  wake placement. It does not need priorities, preemption or scheduler
  callbacks.
\item \textbf{Bounded wait.} No request waits more than a fixed number of
  passes, whatever its usage.
\end{enumerate}
G2 and G3 conflict in one measurable way: fairness shifts the served mix
toward cheap operations, and each operation carries a fixed lock cost. We
quantify this conflict in \S\ref{sec:q5} rather than hide it in an ops/s
number.

% ---------------------------------------------------------------------------
\section{Design}\label{sec:design}

\subsection{Executor contract}
A \lk{Waker} can only enqueue its task, so it carries no placement. The
executor exposes a thread-local placement hint, which the schedule callback
reads and resets:
\begin{itemize}
\item \lk{wake\_with(Inline, w)} puts the task in this worker's run-next slot.
  The slot is polled as soon as the current poll returns, before the local
  queue and before stealing. It has no anti-starvation cap; that is
  deliberate, so that unbounded CES chains remain observable.
\item \lk{wake\_with(Remote, w)} pushes the task to the injector and unparks
  one worker.
\item \lk{wake\_with(Home, w)} pushes the task to the inbox of the worker that
  last polled it.
\item \lk{reschedule\_self\_remote(cx)} marks the \emph{current} task
  \lk{Remote}. \lk{async-task} defers the schedule callback until the poll
  returns \lk{Pending}.
\end{itemize}
Tokio's LIFO slot is an implicit, capped \lk{Inline} for every wake issued on a
worker. The contract makes that choice explicit and per-wake. It adds no
policy to the executor beyond what the lock requests (G4).

\subsection{Burden: CES chain bound and break placement}
A \emph{chain} is the sequence of inline resumes on one worker since that
worker last polled anything else. \lk{ces-k$K$} ends a chain after $K$
hand-offs; \lk{ces-t$T$} ends it after $T$ cycles. At that point the unlocker
does not resume the head waiter inline. It hands ownership to the head waiter
with a \emph{break placement} and continues, so the worker drains its own queue
when the poll returns. With the default break placement, the woken owner lands
in the ex-combiner's own queue and the next chain restarts there. The
\lk{-home} suffix sends it to its home worker's inbox, which moves the
combiner role (\S\ref{sec:q1}).

\subsection{Burden: FC yield, placement and pass limit}
Every FC client owns one node, allocated once. A request writes a closure
pointer and a trampoline into the node and pushes the node onto a Treiber
stack. The client then tries the combiner flag. The winner drains the stack
into the policy queue and runs at most $H$ closures (default 64). Each served
waiter is marked \lk{COMPLETE} (release/acquire) and woken with the lock's
\emph{wake placement}.

Losers return \lk{Pending}; a waiter never spins. At the end of a pass the
combiner handles three cases:
\begin{itemize}
\item Queue empty: release the flag and re-check the stack (SeqCst on both
  sides).
\item Own request served: release the flag and wake the policy's next
  candidate, which re-runs the election.
\item Own request still pending: keep combining.
\end{itemize}
\emph{Yield-after-combine} (default on) makes a poll that won the election
return \lk{Pending} once, having woken itself behind the waiters it just
served. Preemption gives OS-thread FC the same effect. \lk{-remote} and
\lk{-home} set the placement of every wake the lock issues.

\subsection{Service: usage-ordered FC-PQ}
FC-PQ keeps pending requests in a binary min-heap keyed by
(cumulative charged cycles, arrival). The combiner charges each closure the
\lk{rdtscp} cycles around its execution, and the charge persists in the
client's node. Two rules modify the key:
\begin{itemize}
\item A \emph{newcomer}, meaning a client never served before, enters at the
  lock's running mean cost per request. Zero, the minimum and the median are
  alternatives.
\item The \emph{starvation clamp} runs at every pass start. An entry that has
  waited more than $c$ passes (default 8) has its key lowered to the current
  heap minimum until it is served. The accounting is untouched. This fixes a
  no-op in the reference implementation, which clamped after the pop.
\end{itemize}
The pass limit $H$ decides whether ordering matters at all. When
$H\ge$ the number of waiters, every pass admits everyone, ordering only
permutes requests within a pass, and the share stays FIFO's.

\paragraph{Why clamp 8 binds and clamp $\ge$12 does not.}
Take $H{=}16$ with $N_h{=}32$ heavy clients. A heavy client's service interval
is $2(1+\text{L:H})$ passes. Under clamp 8 every heavy operation is a clamp
promotion: \PromotedCEight{} of requests are promoted, which equals the heavy
share of operations $1/(1+3.64)$. The interval is therefore $(c+1)+0.28=9.28$
passes, where 0.28 is the return delay (parallel work plus wake). That gives
L:H$\,=3.64$.

Without the clamp, usage ordering settles where charged usage is equal. That
happens at an interval of $2\times6.51=13.0$ passes. A clamp therefore binds
only if $c+1.28<13.0$, i.e.\ $c\le11$. At clamp 16 it promotes at most
0.035\,\% of requests, all of them in the wait tail.

Service Jain follows from the two-class model with
$x=\text{L:H}\cdot\text{CS}_L/\text{CS}_H$. The model predicts 0.9395 at clamp 8
(measured 0.940) and 0.9970 at clamp 16 (measured 0.997). $J\ge0.95$ needs
$x\ge\XNinetyFive$, i.e.\ L:H\,$\ge\LHNinetyFive$. By the binding rule that
means $c\ge9$. For clamps 9--11 the model predicts $J\approx0.963$, 0.981 and
0.992. We did not run them.

Jain stops at 0.997 rather than 1 because the charge window includes the
trampoline that moves the closure and its result. That is about 193 cycles per
operation that the harness's CS timer does not see.

\paragraph{Utilisation.}
Let $\bar C$ be the mean CS cycles per operation of the served mix and
$o=(T-\mathit{CS})/\mathit{ops}$ the per-operation cycles not spent in a CS.
For FC-family locks under sustained load, where some combiner is busy
98.8--98.9\,\% of the window, $o$ is the lock's per-op cost:
in-pass administration plus hand-off gap. Then lock utilisation is
$\mathit{CS}/T=\bar C/(\bar C+o)$. This is an identity. Its content is
empirical: $o$ is nearly independent of the clamp setting (\OCEight{} cycles
at clamp 8, \OCSixteen{} at clamp 16). Moving the mix toward cheap operations
lowers $\bar C$ (\CbarCEight{} to \CbarCSixteen{} cycles) and therefore lowers
utilisation.

\subsection{The obvious alternative: an actor}
\lk{actor} spawns one server task per lock, lazily with the first request.
Each server poll drains the request stack, serves up to 64 closures in FIFO
order and yields. The server parks when there is nothing to serve.
\lk{actor-inline} changes three placements: the server is woken into the
publisher's run-next slot, it yields \lk{home}, and it wakes clients
\lk{remote}. The actor idiom is what an application would write without a
delegation lock.

% ---------------------------------------------------------------------------
\section{Implementation}\label{sec:impl}

The locks, executor and harness are written in Rust. Table~\ref{tab:loc}
gives the size of each part.

\paragraph{Executor.} Tasks are \lk{async-task} tasks. Each worker has a
\lk{crossbeam-deque} FIFO and a run-next slot; there is one shared
\lk{Injector}. Idle workers park through a SeqCst-fenced registration protocol
that forbids lost wake-ups. Workers are pinned to distinct physical cores.

\paragraph{Requests.} Clients never allocate per request. Each owns one node,
and the lock keeps all nodes alive so that a combiner never touches a freed
node. \lk{Run} futures hold the closure and the result slot inline.
Completion is a release store of \lk{COMPLETE}, and the owner reads it with an
acquire load. All locks share one \lk{rdtscp} counter with the harness.

\paragraph{Checks.} By the lock contract, no lock allocates per request in
steady state: per-client nodes and in-place heaps or queues only. A counting-allocator probe in steady state measured
0~allocations per op for \lk{fc} and \lk{actor}. The executor queues allocate
0.012--0.016 per op for \lk{fc-remote} and 0.016--0.025 for \lk{actor-inline} (one
\lk{Injector} block per 63 remote wakes). \lk{dispatch}, \lk{ces} and \lk{fcpq}
were not probed \todo{allocation probe for \lk{dispatch}/\lk{ces}/\lk{fcpq} not
recorded}.

The \lk{actor} unit tests exercise mutual exclusion with overlap detection, the
absence of lost wake-ups across idle transitions, and FIFO order. They pass in
release builds and under ThreadSanitizer (0 reports). We have no TSan record
for the other locks \todo{TSan runs for \lk{dispatch}/\lk{ces}/\lk{fc}/\lk{fcpq}
not recorded}. \lk{coro-bench --sanity} and \lk{tokio-bench --sanity} (map
length equals total ops) pass for all locks.

\begin{table}[t]
\centering\footnotesize
\adjustbox{max width=\linewidth}{%% generated by scripts/make_tables.py -- do not edit
\begin{tabular}{lrr}
\toprule
component & code & tests \\
\midrule
executor & 681 & 0 \\
lock trait + \texttt{dispatch} & 305 & 0 \\
\texttt{ces} & 180 & 0 \\
\texttt{fc} (delegation core) & 379 & 428 \\
\texttt{fcpq} policy & 281 & 298 \\
\texttt{actor}, \texttt{actor-inline} & 479 & 224 \\
stats, workload, \texttt{coro-bench} & 1119 & 27 \\
tokio baseline (\texttt{tokio-bench}) & 765 & 0 \\
\midrule
total & 4189 & 977 \\
\bottomrule
\end{tabular}
}
\caption{Lines of Rust: non-blank, non-comment lines, with lines from the
  first \lk{\#[cfg(test)]} on counted as tests. Counted by
  \lk{make\_tables.py}.}
\label{tab:loc}
\end{table}

% ---------------------------------------------------------------------------
\section{Evaluation}\label{sec:eval}

\paragraph{Setup.} The machine has 2$\times$ Intel Xeon Gold 6438M (32 cores
per socket, SMT on, 128 logical CPUs), Linux 6.17.7, a 2.20\,GHz TSC and
\lk{rustc} 1.100.0-nightly with \lk{--release}. $W$ workers are pinned to
logical CPUs $0..W{-}1$, one per physical core. Each run has a 200\,ms warm-up
and a 2\,s window. Each configuration is repeated 3 times, with repeats as the
outermost loop, and runs execute one at a time under a measurement lock. We
report medians with [min, max] over the repeats, and we do not interpret
differences inside the spread. Latencies are histogram bucket lower bounds at
6\,\% resolution. One unrelated single-threaded process ran during the 09-28
matrices. \todo{commit id of the measured binaries not recorded; binaries are
identified by sha256 in FINDINGS.md}

\paragraph{Two measurement windows.}
From 2026-09-29 00:07:54 UTC, CPUs 0--15 were capped at 3.0\,GHz. On 09-28
they ran at about 3.69\,GHz turbo. Spins are TSC-timed, so the cap slows only
non-spin work: a light CS costs 1\,365--1\,376 instead of 1\,303 cycles. Every
table row is labelled with its window, and ratios are formed only within one
window unless marked otherwise.

The same binary paths re-run on 09-29 were 1.4--3.4\,\% slower than their
09-28 cells (\lk{dispatch}, \lk{ces-k64-home}, \lk{fc-remote}). This is
outside both spreads, and service Jain rose by 0.004--0.007. We attribute the
drift to the cap [inference; the actor entry that measured it predates the
diagnosis]. Cross-window ratios are therefore uncertain by about 3\,\%.

\subsection{Q1: Does placement restore burden fairness?}\label{sec:q1}

\begin{table*}[t]
\centering\footnotesize
\adjustbox{max width=\linewidth}{%% generated by scripts/make_tables.py -- do not edit
\begin{tabular}{llrrrrrr}
\toprule
cell & variant & burden $J$ & comb.\ p99 & other p99 & starved c/b & Mops/s & /\texttt{ces} \\
\midrule
W8 sus. b0 & \texttt{ces} & 0.125 & cens. & 2.6\rng{2.6}{2.7} & 5\rng{1}{7}/1 & 0.386\rng{0.360}{0.411} & 1.00 \\
W8 sus. b0 & \texttt{ces-k64} & 0.125 & 171.3\rng{171.3}{178.7} & 2.7 & 0/0 & 0.377\rng{0.377}{0.379} & 0.98 \\
W8 sus. b0 & \texttt{ces-k64-home} & 0.996\rng{0.993}{0.997} & 2.6\rng{2.6}{2.8} & 2.6\rng{2.6}{2.8} & 0/0 & 0.378\rng{0.378}{0.382} & 0.98 \\
W8 sus. b0 & \texttt{ces-t64000-home} & 0.999\rng{0.999}{1.000} & 2.6 & 2.6 & 0/0 & 0.354\rng{0.354}{0.356} & 0.92 \\
W8 sus. b0 & \texttt{fc} & 0.125 & 268.1\rng{268.1}{283.0} & 0.1 & 0/0 & 0.234\rng{0.232}{0.235} & 0.61 \\
W8 sus. b0 & \texttt{fc-home} & 0.744\rng{0.725}{0.829} & 2.3\rng{2.2}{2.4} & 2.4 & 0/0 & 0.369\rng{0.367}{0.373} & 0.96 \\
W8 sus. b0 & \texttt{fc-remote} & 0.981\rng{0.967}{0.996} & 2.6\rng{2.6}{2.7} & 2.6\rng{2.6}{2.7} & 0/0 & 0.388\rng{0.382}{0.388} & 1.01 \\
\midrule
W8 sus. b31 & \texttt{ces} & 0.125 & -- & 3.0 & 0/0 & 0.385\rng{0.384}{0.387} & 1.00 \\
W8 sus. b31 & \texttt{ces-k64} & 0.904\rng{0.897}{0.978} & 3.0 & 3.0 & 0/0 & 0.380\rng{0.379}{0.380} & 0.99 \\
W8 sus. b31 & \texttt{ces-k64-home} & 0.993\rng{0.992}{0.995} & 2.7\rng{2.7}{2.8} & 2.9\rng{2.8}{3.0} & 0/0 & 0.378\rng{0.377}{0.379} & 0.98 \\
W8 sus. b31 & \texttt{ces-t64000-home} & 1.000\rng{0.999}{1.000} & 2.6 & 2.6 & 0/0 & 0.355\rng{0.354}{0.355} & 0.92 \\
W8 sus. b31 & \texttt{fc} & 0.999\rng{0.998}{0.999} & 1.3\rng{1.2}{1.3} & 1.3 & 0/0 & 0.392\rng{0.387}{0.393} & 1.02 \\
W8 sus. b31 & \texttt{fc-home} & 0.986\rng{0.975}{0.994} & 2.6\rng{2.4}{2.6} & 2.6\rng{2.6}{2.7} & 0/0 & 0.378\rng{0.377}{0.378} & 0.98 \\
W8 sus. b31 & \texttt{fc-remote} & 0.993\rng{0.975}{0.995} & 2.9\rng{2.9}{3.0} & 3.0 & 0/0 & 0.387\rng{0.386}{0.387} & 1.00 \\
\midrule
W8 bur. b0 & \texttt{ces} & 0.999\rng{0.999}{1.000} & 44.7 & 44.7 & 0/0 & 0.317\rng{0.317}{0.318} & 1.00 \\
W8 bur. b0 & \texttt{ces-k64} & 0.999\rng{0.999}{1.000} & 44.7 & 44.7 & 0/0 & 0.317\rng{0.316}{0.317} & 1.00 \\
W8 bur. b0 & \texttt{ces-k64-home} & 0.996\rng{0.993}{0.999} & 44.7 & 44.7 & 0/0 & 0.314\rng{0.313}{0.317} & 0.99 \\
W8 bur. b0 & \texttt{ces-t64000-home} & 1.000\rng{0.998}{1.000} & 29.8 & 29.8 & 0/0 & 0.300\rng{0.292}{0.301} & 0.95 \\
W8 bur. b0 & \texttt{fc} & 0.125 & 268.1 & 0.1 & 0/0 & 0.059 & 0.19 \\
W8 bur. b0 & \texttt{fc-home} & 0.946\rng{0.946}{0.947} & 37.2 & 37.2 & 0/0 & 0.358\rng{0.357}{0.358} & 1.13 \\
W8 bur. b0 & \texttt{fc-remote} & 1.000 & 29.8 & 29.8 & 0/0 & 0.304\rng{0.301}{0.304} & 0.96 \\
\midrule
W8 bur. b31 & \texttt{ces} & 1.000 & 44.7 & 44.7 & 0/0 & 0.317\rng{0.316}{0.317} & 1.00 \\
W8 bur. b31 & \texttt{ces-k64} & 1.000 & 44.7 & 44.7 & 0/0 & 0.317\rng{0.317}{0.318} & 1.00 \\
W8 bur. b31 & \texttt{ces-k64-home} & 0.996\rng{0.995}{0.998} & 44.7 & 44.7 & 0/0 & 0.314\rng{0.314}{0.316} & 0.99 \\
W8 bur. b31 & \texttt{ces-t64000-home} & 1.000 & 29.8 & 29.8 & 0/0 & 0.306\rng{0.299}{0.306} & 0.97 \\
W8 bur. b31 & \texttt{fc} & 1.000 & 29.8 & 29.8 & 0/0 & 0.302\rng{0.300}{0.302} & 0.95 \\
W8 bur. b31 & \texttt{fc-home} & 1.000 & 26.1\rng{26.1}{27.0} & 27.0\rng{27.0}{27.9} & 0/0 & 0.346\rng{0.340}{0.346} & 1.09 \\
W8 bur. b31 & \texttt{fc-remote} & 1.000 & 29.8 & 29.8 & 0/0 & 0.307\rng{0.306}{0.307} & 0.97 \\
\midrule
W16 sus. b0 & \texttt{ces} & 0.062 & cens. & 2.7 & 2\rng{2}{3}/1 & 0.375\rng{0.371}{0.390} & 1.00 \\
W16 sus. b0 & \texttt{ces-k64-home} & 0.999\rng{0.996}{0.999} & 2.7 & 2.7 & 0/0 & 0.372\rng{0.372}{0.374} & 0.99 \\
\midrule
W16 sus. b31 & \texttt{ces} & 0.062 & -- & 2.6 & 0/0 & 0.377\rng{0.376}{0.380} & 1.00 \\
W16 sus. b31 & \texttt{ces-k64-home} & 0.998\rng{0.995}{0.998} & 2.6 & 2.6 & 0/0 & 0.369\rng{0.369}{0.372} & 0.98 \\
\midrule
W16 bur. b0 & \texttt{ces} & 0.062\rng{0.062}{0.096} & cens. & 14.9 & 0/1 & 0.365\rng{0.363}{0.366} & 1.00 \\
W16 bur. b0 & \texttt{ces-k64-home} & 0.999\rng{0.998}{0.999} & 14.9 & 14.9 & 0/0 & 0.346\rng{0.345}{0.346} & 0.95 \\
\midrule
W16 bur. b31 & \texttt{ces} & 0.062 & -- & 14.9 & 0/0 & 0.364\rng{0.361}{0.366} & 1.00 \\
W16 bur. b31 & \texttt{ces-k64-home} & 0.999\rng{0.998}{0.999} & 14.9 & 14.9 & 0/0 & 0.348 & 0.96 \\
\bottomrule
\end{tabular}
}
\caption{Burden fairness, phase 3, heavy $8\times$ (09-28 window). The p99
  columns are bystander delays in $\mu$s. ``/\lk{ces}'' is the throughput ratio
  to plain CES in the same cell. \lk{ces} rows with ``--'' combiner p99 have no
  bystander left on the combiner worker (evicted by balancing).}
\label{tab:burden}
\end{table*}

Table~\ref{tab:burden} gives the answer.

\emph{A chain bound alone does nothing for burden at b0.}
\lk{ces-k64} keeps burden at 0.125. The ex-combiner drains its own queue, the
next acquirer is again one of its own clients, and the combiner's bystander
p99 is 171\,$\mu$s, the length of one 64-hand-off chain. The bound ends
starvation but not concentration.

\emph{The bound plus a home chain break fixes it.}
\lk{ces-k64-home} has burden Jain 0.993--0.999 in every cell (8 and 16
workers, sustained and bursty, b0 and b31). Its combiner-worker bystander p99
equals the others' (2.6--2.9\,$\mu$s sustained; 44.7 and 14.9\,$\mu$s bursty
at 8 and 16 workers), and no task starves. It costs 0.98--0.99$\times$ CES
throughput sustained and 0.95--0.99$\times$ bursty.

\emph{Bursty load needs no bound.} In bursty cells CES chains already end
naturally (p50 12 hand-offs). The bound only trims the tail from a maximum of
152--191 to 64. The cycle budget (\lk{ces-t64000-home}) breaks every
${\sim}12$ hand-offs even where the queue would have continued, and costs
5--8\,\%.

\emph{For FC, \lk{remote} beats \lk{home}.} \lk{fc-remote} removes the b0
one-worker convergence (burden 0.981\rng{0.967}{0.996} sustained, 1.000
bursty) at the balanced throughput level. \lk{fc-home} reaches only 0.744 at
b0 sustained. The election is sticky: the ex-combiner's own clients are woken
locally and win the next \lk{try\_lock}. In bursty mode \lk{home} is
nevertheless faster (1.09--1.13$\times$ CES), because served waiters resume
their parallel work on their own worker.

\emph{Bystander delay did not exceed the other workers' under balancing.}
Our pre-registered hypothesis was that the combiner-worker bystander p99 would
exceed the other workers' p99 by more than one pass ($H\times$ mean CS). It
is refuted under balancing in all 12 bursty cells for all five combining
variants: the two p99s are equal within one bucket, for example 44.7 against
44.7\,$\mu$s for CES at 8 workers, with a 136.6\,$\mu$s pass. Without balancing
the effect is starvation, not delay.

\subsection{Q2: Does usage ordering restore service fairness?}\label{sec:q2}

\begin{table*}[t]
\centering\footnotesize
\adjustbox{max width=\linewidth}{%% generated by scripts/make_tables.py -- do not edit
\begin{tabular}{llrrrrrrrr}
\toprule
variant & win. & Mops/s & svc.\ $J$ & model $J$ & L:H & util & non-CS/op & heavy p50 / p99 ($\mu$s) & max wait \\
\midrule
\texttt{dispatch} & 09-28 & 0.217 & 0.659\rng{0.659}{0.660} & 0.659 & 1.00 & 0.482 & 5249 & 283.0 / 297.9 & -- \\
\texttt{ces} & 09-28 & 0.385\rng{0.384}{0.387} & 0.651\rng{0.651}{0.652} & 0.651 & 1.00 & 0.838 & 924 & 156.4 / 178.7 & -- \\
\texttt{fc} & 09-28 & 0.392\rng{0.387}{0.393} & 0.655\rng{0.654}{0.658} & 0.655 & 1.00 & 0.860 & 784 & 156.4 / 238.3 & -- \\
\texttt{fc-remote} & 09-28 & 0.387\rng{0.386}{0.387} & 0.654 & 0.654 & 1.00 & 0.848 & 862 & 156.4 / 178.7\rng{171.3}{178.7} & -- \\
\texttt{fcpq} & 09-28 & 0.430\rng{0.429}{0.431} & 0.706\rng{0.706}{0.714} & 0.706 & 1.36 & 0.838 & 828 & 163.8 / 268.1 & -- \\
\texttt{fcpq-h8} & 09-28 & 0.464\rng{0.463}{0.465} & 0.763\rng{0.763}{0.764} & 0.764 & 1.80 & 0.806 & 921 & 186.2 / 268.1 & -- \\
\texttt{fcpq-h16} & 09-28 & 0.547\rng{0.546}{0.547} & 0.863\rng{0.862}{0.863} & 0.863 & 2.73 & 0.792 & 838 & 216.0 / 312.8 & -- \\
\texttt{fcpq-h16-home} & 09-28 & 0.564\rng{0.563}{0.571} & 0.932\rng{0.930}{0.933} & 0.932 & 3.65 & 0.725 & 1071 & 253.2 / 342.5\rng{342.5}{342.6} & -- \\
\midrule
\texttt{fc-remote} & 09-29 & 0.378\rng{0.377}{0.379} & 0.660 & 0.660 & 1.00 & 0.839 & 937 & 163.8 / 178.7 & -- \\
\texttt{fcpq-h16-home} & 09-29 & 0.542\rng{0.541}{0.543} & 0.939 & 0.939 & 3.64 & 0.710 & 1175 & 253.2 / 357.4 & -- \\
\quad\texttt{-c8} & 09-29 & 0.540\rng{0.539}{0.540} & 0.940\rng{0.939}{0.940} & 0.940 & 3.64 & 0.708 & 1189 & 253.2 / 357.4 & 10 \\
\quad\texttt{-c16} & 09-29 & 0.612\rng{0.611}{0.614} & 0.997 & 0.997 & 5.51 & 0.681 & 1145 & 327.7 / 625.5 & 17 \\
\quad\texttt{-c32} & 09-29 & 0.613\rng{0.612}{0.613} & 0.997 & 0.997 & 5.51 & 0.682 & 1140 & 327.7 / 625.5 & 33 \\
\quad\texttt{-c0} & 09-29 & 0.611\rng{0.610}{0.612} & 0.997 & 0.997 & 5.48 & 0.681 & 1147 & 327.7 / 625.5 & 106\rng{54}{193} \\
\bottomrule
\end{tabular}
}
\caption{Service fairness, $W{=}8$, sustained, b31. Upper block: phase 3
  (09-28). Lower block: same-window A/B and clamp sweep (09-29, 3.0\,GHz cap);
  \lk{-c$N$} is \lk{fcpq-h16-home} with starvation clamp $N$ (0 = off) and
  newcomer init \emph{mean}. ``model $J$'' is the two-class model applied to
  the measured L:H and per-class costs. L:H, util and non-CS/op are medians (max spread 5\,\% of the median). util $=$ CS cycles / window. non-CS/op $=$
  (window $-$ CS cycles) / ops; only for the FC-family rows, whose combiner is busy almost the whole window, is it the per-op lock cost $o$ of \S\ref{sec:design}; for \lk{dispatch} and \lk{ces} it also contains lock-idle and hand-off time. Max wait is in combining passes. Do not compare
  rows across the two blocks.}
\label{tab:service}
\end{table*}

\begin{table}[t]
\centering\footnotesize\setlength{\tabcolsep}{3pt}
\adjustbox{max width=\linewidth}{%% generated by scripts/make_tables.py -- do not edit
\begin{tabular}{llrrrrr}
\toprule
cell & clamp & Mops/s & svc.\ $J$ & L:H & heavy p99 ($\mu$s) & burden $J$ \\
\midrule
W8 sus. b31 & 8 & 0.540\rng{0.539}{0.540} & 0.940\rng{0.939}{0.940} & 3.64 & 357.4 & 1.000 \\
W8 sus. b31 & 16 & 0.612\rng{0.611}{0.614} & 0.997 & 5.51 & 625.5 & 1.000 \\
\midrule
W8 sus. b0 & 8 & 0.540\rng{0.540}{0.543} & 0.930\rng{0.929}{0.931} & 3.50 & 268.1 & 0.501 \\
W8 sus. b0 & 16 & 0.614\rng{0.613}{0.617} & 0.997 & 5.53 & 655.3 & 0.970 \\
\midrule
W16 sus. b31 & 8 & 0.531\rng{0.530}{0.531} & 0.937\rng{0.937}{0.938} & 3.57 & 342.5 & 0.999 \\
W16 sus. b31 & 16 & 0.604\rng{0.603}{0.604} & 0.997 & 5.45 & 565.9 & 0.998 \\
\midrule
W8 bur. b31 & 8 & 0.358\rng{0.357}{0.358} & 0.697\rng{0.696}{0.697} & 1.23 & 78.2 & 1.000 \\
W8 bur. b31 & 16 & 0.358\rng{0.357}{0.358} & 0.697\rng{0.696}{0.697} & 1.23 & 78.2 & 1.000 \\
\midrule
W16 bur. b31 & 8 & 0.386 & 0.734\rng{0.734}{0.735} & 1.48 & 70.7 & 0.998 \\
W16 bur. b31 & 16 & 0.386 & 0.734\rng{0.734}{0.735} & 1.48 & 70.7 & 0.998 \\
\bottomrule
\end{tabular}
}
\caption{Clamp 8 vs.\ 16 for \lk{fcpq-h16-home} in the confirmation cells
  (09-29 window). Heavy p99 is run latency in $\mu$s; L:H, p99 and burden $J$ are medians.}
\label{tab:confirm}
\end{table}

The pass limit decides whether usage ordering has any effect
(Table~\ref{tab:service}). \lk{fcpq} with $H{=}64$ admits every waiter in every
pass, so it reaches only 0.706. $H{=}8$ gives 0.763, $H{=}16$ 0.863, and
$H{=}16$ with \lk{home} wakes 0.932. The remaining FC-PQ knobs leave service Jain at or below \lk{fcpq}'s level:
pass budget \lk{t16000} 0.661, rotation 0.708, combining credit 0.693,
max-usage election 0.707. None of them changes the admission set.

At $H{=}16$ with \lk{home} wakes, the 8-pass clamp is the bound: L:H stays at
3.64, as \S\ref{sec:design} predicts. Lengthening the clamp to 16 lifts
service Jain to 0.997 at $W{=}8$ b31, $W{=}8$ b0 and $W{=}16$ b31
(Table~\ref{tab:confirm}). L:H rises to 5.45--5.53, and ops throughput rises
by 13\,\% (0.540$\to$0.612\,Mops/s, same window).

The price is heavy-client latency. Heavy p50/p99 grow from 253/357\,$\mu$s to
328/626\,$\mu$s, and light p99 from 164 to 343\,$\mu$s. The worst queue wait
is 17 passes at clamp 16 and 33 at clamp 32. With the clamp off it is
106\rng{54}{193} passes, with heavy maximum latencies of 1.5--4.3\,ms.
We therefore take \lk{fcpq-h16-home-c16}: it has the fairness of clamp-off and
a bounded worst case (G5).

Two settings had no effect. The newcomer rule fires on only 64 of about
1.2\,M requests per run, all during warm-up when every usage is near 0, so
mean, zero, min and median initialisation are indistinguishable. In bursty
load every knob is inert (Jain 0.697 at $W{=}8$): with at most 16 waiters and
7.8 ops per pass, every pass serves everyone pending. Usage ordering can only
act when the backlog exceeds $H$.

\emph{Side effect on burden.} Clamp 16 also fixes a burden problem at b0, from
0.501\rng{0.500}{0.544} to 0.970\rng{0.969}{0.974}. Our explanation
[inference]: under clamp 8 the pass-end hand-off almost always goes to a light
client, and at b0 light clients live on 4 of the 8 workers.

\subsection{Q3: Against real tokio locks}\label{sec:q3}

\begin{table*}[t]
\centering\footnotesize\setlength{\tabcolsep}{3.5pt}
\adjustbox{max width=\linewidth}{%% generated by scripts/make_tables.py -- do not edit
\begin{tabular}{lrrrrrrrr}
\toprule
 & \multicolumn{2}{c}{W8 sus.} & \multicolumn{2}{c}{W8 bur.} & \multicolumn{2}{c}{W16 sus.} & \multicolumn{2}{c}{W16 bur.} \\
\cmidrule(lr){2-3} \cmidrule(lr){4-5} \cmidrule(lr){6-7} \cmidrule(lr){8-9}
variant & Mops/s & $\times$tm & Mops/s & $\times$tm & Mops/s & $\times$tm & Mops/s & $\times$tm \\
\midrule
\texttt{tokio-mutex} & 0.241\rng{0.241}{0.242} & 1.00 & 0.073 & 1.00 & 0.232\rng{0.231}{0.233} & 1.00 & 0.072 & 1.00 \\
\texttt{tokio-mutex-unconstrained} & 0.239 & 0.99 & 0.073 & 1.00 & 0.231\rng{0.230}{0.231} & 0.99 & 0.072 & 1.00 \\
\texttt{async-lock} & 0.177\rng{0.177}{0.239} & 0.73 & 0.054 & 0.75 & 0.176 & 0.76 & 0.054 & 0.75 \\
\texttt{std-mutex} & 0.204\rng{0.203}{0.230} & 0.85 & 0.366\rng{0.365}{0.384} & 5.02 & 0.167\rng{0.166}{0.167} & 0.72 & 0.195\rng{0.172}{0.228} & 2.70 \\
\texttt{parking-lot} & 0.140\rng{0.138}{0.142} & 0.58 & 0.357\rng{0.356}{0.359} & 4.90 & 0.108\rng{0.108}{0.114} & 0.46 & 0.125\rng{0.122}{0.154} & 1.73 \\
\midrule
\texttt{dispatch} & 0.218\rng{0.217}{0.218} & 0.90 & 0.071 & 0.97 & 0.214\rng{0.214}{0.215} & 0.92 & 0.071 & 0.98 \\
\texttt{dispatch-home} & 0.253 & 1.05 & 0.147\rng{0.146}{0.150} & 2.01 & 0.239\rng{0.239}{0.241} & 1.03 & 0.172\rng{0.161}{0.172} & 2.39 \\
\texttt{ces-k64-home} & 0.381\rng{0.380}{0.381} & 1.58 & 0.317\rng{0.314}{0.318} & 4.36 & 0.372 & 1.60 & 0.350\rng{0.349}{0.350} & 4.85 \\
\texttt{fc-remote} & 0.387\rng{0.387}{0.388} & 1.61 & 0.307\rng{0.307}{0.308} & 4.22 & 0.378\rng{0.378}{0.379} & 1.63 & 0.358\rng{0.357}{0.358} & 4.97 \\
\texttt{fcpq-h16-home} & 0.571\rng{0.569}{0.571} & 2.37 & 0.365\rng{0.365}{0.366} & 5.01 & 0.547\rng{0.545}{0.549} & 2.36 & 0.392 & 5.44 \\
\bottomrule
\end{tabular}
}
\caption{Throughput against tokio 1.53.1 (09-28 window; coro rows b31).
  $\times$tm is the ratio to \lk{tokio-mutex} in the same cell.
  \lk{fcpq-h16-home} (clamp 8) counts cheap light ops and is not a
  like-for-like throughput reference. \textbf{The bursty ratios of 4--5$\times$
  are largely due to tokio's LIFO slot; see Table~\ref{tab:lifo}.}}
\label{tab:xrt}
\end{table*}

\begin{table*}[t]
\centering\footnotesize\setlength{\tabcolsep}{3.5pt}
\adjustbox{max width=\linewidth}{%% generated by scripts/make_tables.py -- do not edit
\begin{tabular}{lrrrrrrrr}
\toprule
 & \multicolumn{2}{c}{W8 sus.} & \multicolumn{2}{c}{W8 bur.} & \multicolumn{2}{c}{W16 sus.} & \multicolumn{2}{c}{W16 bur.} \\
\cmidrule(lr){2-3} \cmidrule(lr){4-5} \cmidrule(lr){6-7} \cmidrule(lr){8-9}
variant & svc.\ $J$ & st.\ c/b & svc.\ $J$ & st.\ c/b & svc.\ $J$ & st.\ c/b & svc.\ $J$ & st.\ c/b \\
\midrule
\texttt{tokio-mutex} & 0.663\rng{0.662}{0.663} & 0/0 & 0.668\rng{0.668}{0.669} & 0/0 & 0.666\rng{0.665}{0.668} & 0/0 & 0.672\rng{0.671}{0.674} & 0/0 \\
\texttt{tokio-mutex-unconstrained} & 0.665\rng{0.664}{0.665} & 0/0 & 0.669\rng{0.669}{0.670} & 0/0 & 0.667 & 0/0 & 0.672 & 0/0 \\
\texttt{async-lock} & 0.016\rng{0.016}{0.661} & 63\rng{0}{63}/0\rng{0}{2} & 0.062 & 15/0\rng{0}{1} & 0.016 & 63/0 & 0.062 & 15/0 \\
\texttt{std-mutex} & 0.055\rng{0.047}{0.057} & 56/8 & 0.255\rng{0.253}{0.256} & 8/8 & 0.234 & 48/16 & 0.782\rng{0.767}{0.807} & 0/16 \\
\texttt{parking-lot} & 0.048\rng{0.048}{0.049} & 56/8 & 0.250 & 8/8 & 0.197\rng{0.181}{0.198} & 48/16 & 0.774\rng{0.692}{0.775} & 0/16 \\
\midrule
\texttt{dispatch} & 0.659 & 0/0 & 0.674 & 0/0 & 0.663 & 0/0 & 0.679 & 0/0 \\
\texttt{dispatch-home} & 0.669 & 0/0 & 0.669\rng{0.669}{0.670} & 0/0 & 0.674\rng{0.673}{0.674} & 0/0 & 0.673\rng{0.672}{0.673} & 0/0 \\
\texttt{ces-k64-home} & 0.659 & 0/0 & 0.678\rng{0.677}{0.679} & 0/0 & 0.661\rng{0.661}{0.662} & 0/0 & 0.663 & 0/0 \\
\texttt{fc-remote} & 0.653 & 0/0 & 0.655\rng{0.655}{0.656} & 0/0 & 0.659\rng{0.659}{0.660} & 0/0 & 0.662 & 0/0 \\
\texttt{fcpq-h16-home} & 0.930 & 0/0 & 0.685\rng{0.685}{0.686} & 0/0 & 0.933\rng{0.932}{0.934} & 0/0 & 0.726\rng{0.724}{0.727} & 0/0 \\
\bottomrule
\end{tabular}
}
\caption{Service Jain and starved clients/bystanders for the runs of
  Table~\ref{tab:xrt}. \lk{std-mutex} and \lk{parking-lot} are really $W$
  pinned threads on a blocking mutex (see text).}
\label{tab:xrtfair}
\end{table*}

\begin{table*}[t]
\centering\footnotesize\setlength{\tabcolsep}{2.5pt}
\adjustbox{max width=\linewidth}{%% generated by scripts/make_tables.py -- do not edit
\begin{tabular}{lrrrrrrr}
\toprule
 & \multicolumn{2}{c}{\texttt{tokio-mutex}} & \multicolumn{2}{c}{\texttt{async-lock}} & \multicolumn{3}{c}{coro / \texttt{tokio-mutex} LIFO off} \\
\cmidrule(lr){2-3} \cmidrule(lr){4-5} \cmidrule(lr){6-8}
cell & LIFO on & off & on & off & \texttt{ces-k64-home} & \texttt{fc-remote} & \texttt{fcpq-h16-home} \\
\midrule
W8 sus. & 0.231\rng{0.230}{0.242} & 0.250\rng{0.243}{0.251} & 0.176\rng{0.176}{0.403} & 0.237\rng{0.233}{0.241} & 1.52 & 1.55 & 2.28 \\
W8 bur. & 0.072\rng{0.072}{0.073} & 0.263\rng{0.261}{0.265} & 0.072\rng{0.054}{0.072} & 0.250\rng{0.249}{0.253} & 1.21 & 1.17 & 1.39 \\
W16 sus. & 0.229\rng{0.229}{0.232} & 0.248\rng{0.245}{0.253} & 0.176 & 0.234\rng{0.232}{0.244} & 1.50 & 1.53 & 2.21 \\
W16 bur. & 0.072 & 0.248\rng{0.246}{0.251} & 0.054\rng{0.054}{0.072} & 0.237\rng{0.235}{0.244} & 1.41 & 1.44 & 1.58 \\
\bottomrule
\end{tabular}
}
\caption{tokio LIFO slot on/off (Mops/s; throwaway \lk{tokio\_unstable}
  build whose LIFO-on runs are 0.96--1.00$\times$ the release binary; 09-28
  window). The right-hand columns divide the coro throughputs of
  Table~\ref{tab:xrt} by LIFO-off \lk{tokio-mutex}.}
\label{tab:lifo}
\end{table*}

We ported the workload unchanged to tokio's multi-thread runtime
(\lk{tokio-bench}): same key stream, costs, histogram and pinning.

\paragraph{The executor is not a weak baseline.} Coro \lk{dispatch} runs at
0.90--0.98$\times$ \lk{tokio-mutex} (Table~\ref{tab:xrt}). The coop budget
never forced a yield, and removing it (\lk{-unconstrained}) moves throughput
by 0--1\,\%.

\paragraph{Sustained load.} The burden-fair delegation locks run
\SusOnLo--\SusOnHi$\times$ \lk{tokio-mutex}. They starve no client and no
bystander and keep FIFO-level service Jain (0.653--0.678, against 0.663--0.672
for \lk{tokio-mutex}).

\paragraph{The LIFO-slot caveat.} Bursty ratios of
\BurOnLo--\BurOnHi$\times$ come mostly from tokio's LIFO slot, not from the
lock. A mutex hand-off wake lands in the unlocker's LIFO slot, and the new
owner waits there behind the unlocker's parallel work. Disabling the slot
lifts \lk{tokio-mutex} 3.4--3.7$\times$ in bursty mode and 1.08$\times$ in
sustained mode (Table~\ref{tab:lifo}). Against that configuration,
\lk{ces-k64-home} and \lk{fc-remote} run \BurOffLo--\BurOffHi$\times$ bursty
and \SusOffLo--\SusOffHi$\times$ sustained. The advantage that holds under
either tokio configuration is 1.5--1.6$\times$ sustained and about
1.2--1.4$\times$ bursty.

\paragraph{Blocking mutexes.} \lk{std-mutex} and \lk{parking-lot} beat every
async lock at $W{=}8$ bursty (5.02$\times$ and 4.90$\times$). In that cell they
serve only 8 of the 16 clients and starve all bystanders
(Table~\ref{tab:xrtfair}), so the ratio is not like-for-like. Elsewhere the
delegation locks beat \lk{std-mutex} 1.80--3.28$\times$.

\paragraph{Missing comparison.} The two runtimes were only measured together
in the 09-28 window. \lk{fcpq-h16-home-c16} (09-29) has no same-window tokio
reference \todo{same-window tokio baseline for \lk{fcpq-h16-home-c16}}.

\subsection{Q4: Is a server task enough?}\label{sec:q4}

\begin{table*}[t]
\centering\footnotesize
\adjustbox{max width=\linewidth}{%% generated by scripts/make_tables.py -- do not edit
\begin{tabular}{lllrrrrr}
\toprule
cell & variant & win. & Mops/s & /\texttt{fc-remote} & burden $J$ & comb.\ / other p99 ($\mu$s) & svc.\ $J$ \\
\midrule
W8 sus. b31 & \texttt{dispatch} & 09-29 & 0.213\rng{0.212}{0.213} & 0.56 & -- & -- / 0.8 & 0.666\rng{0.665}{0.666} \\
W8 sus. b31 & \texttt{ces-k64-home} & 09-29 & 0.365\rng{0.364}{0.366} & 0.97 & 0.999\rng{0.998}{0.999} & 2.9 / 2.9 & 0.664 \\
W8 sus. b31 & \texttt{fc-remote} & 09-29 & 0.377\rng{0.377}{0.378} & 1.00 & 0.989\rng{0.979}{0.991} & 3.1 / 3.1 & 0.660 \\
W8 sus. b31 & \texttt{actor} & 09-29 & 0.343\rng{0.342}{0.347} & 0.91 & 0.990\rng{0.987}{0.993} & 1.3 / 1.3 & 0.660 \\
W8 sus. b31 & \texttt{actor-inline} & 09-29 & 0.391\rng{0.390}{0.391} & 1.04 & 0.179\rng{0.125}{0.267} & 3.1 / 3.1 & 0.654\rng{0.654}{0.655} \\
\midrule
W8 bur. b31 & \texttt{dispatch} & 09-29 & 0.069 & 0.23 & -- & -- / 0.9 & 0.682\rng{0.681}{0.682} \\
W8 bur. b31 & \texttt{ces-k64-home} & 09-29 & 0.309\rng{0.305}{0.309} & 1.03 & 1.000\rng{0.988}{1.000} & 44.7 / 44.7 & 0.682\rng{0.681}{0.682} \\
W8 bur. b31 & \texttt{fc-remote} & 09-29 & 0.299 & 1.00 & 1.000 & 29.8 / 29.8 & 0.662 \\
W8 bur. b31 & \texttt{actor} & 09-29 & 0.163\rng{0.159}{0.164} & 0.54 & 0.999 & 2.1 / 2.1 & 0.666 \\
W8 bur. b31 & \texttt{actor-inline} & 09-29 & 0.287\rng{0.281}{0.288} & 0.96 & 1.000\rng{0.998}{1.000} & 29.8 / 29.8 & 0.663 \\
\midrule
W8 sus. b0 & \texttt{fc-remote} & 09-28 & 0.388\rng{0.382}{0.388} & 1.00 & 0.981\rng{0.967}{0.996} & 2.6 / 2.6 & 0.654\rng{0.654}{0.658} \\
W8 sus. b0 & \texttt{actor} & 09-29 & 0.228\rng{0.228}{0.231} & 0.59$^\dagger$ & 0.125 & 283.0 / 0.2 & 0.655 \\
W8 sus. b0 & \texttt{actor-inline} & 09-29 & 0.389\rng{0.387}{0.390} & 1.00$^\dagger$ & 0.125 & 156.4 / 2.7 & 0.655\rng{0.654}{0.655} \\
\midrule
W8 bur. b0 & \texttt{fc-remote} & 09-28 & 0.304\rng{0.301}{0.304} & 1.00 & 1.000 & 29.8 / 29.8 & 0.657\rng{0.656}{0.660} \\
W8 bur. b0 & \texttt{actor} & 09-29 & 0.058 & 0.19$^\dagger$ & 0.125 & 268.1 / 0.1 & 0.658 \\
W8 bur. b0 & \texttt{actor-inline} & 09-29 & 0.287\rng{0.283}{0.288} & 0.95$^\dagger$ & 0.999\rng{0.989}{1.000} & 29.8 / 29.8 & 0.663\rng{0.663}{0.664} \\
\midrule
W16 sus. b31 & \texttt{fc-remote} & 09-28 & 0.378\rng{0.378}{0.379} & 1.00 & 0.998\rng{0.994}{0.998} & 2.4 / 2.4 & 0.659\rng{0.659}{0.660} \\
W16 sus. b31 & \texttt{actor} & 09-29 & 0.340 & 0.90$^\dagger$ & 0.970\rng{0.966}{0.970} & 0.9 / 1.0 & 0.661 \\
W16 sus. b31 & \texttt{actor-inline} & 09-29 & 0.382\rng{0.382}{0.383} & 1.01$^\dagger$ & 0.612\rng{0.556}{0.651} & 2.6 / 2.6 & 0.656\rng{0.656}{0.657} \\
\midrule
W16 bur. b31 & \texttt{fc-remote} & 09-28 & 0.358\rng{0.357}{0.358} & 1.00 & 1.000 & 14.9 / 14.9 & 0.662 \\
W16 bur. b31 & \texttt{actor} & 09-29 & 0.173 & 0.48$^\dagger$ & 0.999\rng{0.997}{0.999} & 1.2 / 1.2 & 0.668\rng{0.667}{0.668} \\
W16 bur. b31 & \texttt{actor-inline} & 09-29 & 0.367\rng{0.366}{0.368} & 1.03$^\dagger$ & 0.989\rng{0.988}{0.993} & 14.9 / 14.9 & 0.663 \\
\bottomrule
\end{tabular}
}
\caption{Actor control. Upper blocks: same-window (09-29) references.
  Lower blocks: b0 and $W{=}16$, where the \lk{fc-remote} reference is a
  stored 09-28 cell. Ratios marked $\dagger$ cross windows (actor 09-29 vs.\ reference 09-28) and understate the actor rows
  by up to about 3\,\%.}
\label{tab:actor}
\end{table*}

Table~\ref{tab:actor} compares the two actor variants with the delegation
locks. Plain \lk{actor} is burden-fair only because the executor's balancing
steal keeps moving the server: burden Jain is 0.990 sustained and 0.999 bursty
at $W{=}8$ b31. It runs 0.91$\times$ same-window \lk{fc-remote} sustained and
0.54$\times$ bursty. Served clients are woken into the server's queue ahead of
the yielded server, so each pass waits for their parallel work [inference].
At b0, \lk{actor} is a one-worker system: burden 0.125, 0.59$\times$ sustained,
0.19$\times$ bursty, and a server-worker bystander p99 of 268--283\,$\mu$s.

\lk{actor-inline} restores throughput (1.04$\times$ sustained, 0.96$\times$
bursty at $W{=}8$ b31), but under sustained load the server never parks. A
server that never parks is never re-placed by a wake, so only balancing steals
move it. Burden is 0.179\rng{0.125}{0.267} at $W{=}8$ and
0.612\rng{0.556}{0.651} at $W{=}16$. At b0, its worker's bystander p99 is
156\,$\mu$s, about one 64-request pass.

No actor variant reaches both delegation-level throughput and burden
$\ge0.9$ in the sustained cells. \lk{ces-k64-home} and \lk{fc-remote} reach
both, because they move the combiner role by construction. A server task would
need the same explicit migration.

\subsection{Q5: Costs and limits}\label{sec:q5}

\paragraph{Utilisation.} Service fairness lowers the fraction of time the
lock spends doing work. On 09-28, utilisation drops from \UtilFcOld{}
(\lk{fc}) to 0.725 (\lk{fcpq-h16-home}). On 09-29 it drops from
\UtilFcRemoteAB{} (\lk{fc-remote}) to \UtilCEight{} (clamp 8) and
\UtilCSixteen{} (clamp 16), while a combiner is busy 98.8--98.9\,\% of the
window. The lock's cost is $o=\OCSixteen$ cycles/op, 96--97\,\% of it in-pass
administration. Same-window \lk{fc-remote} has $o=\OFcRemoteAB$, and the
lowest $o$ of any fc-family sustained cell is \lk{fc} on 09-28 at turbo clock,
with $o=\OFcOld$.

At the fairness bound $J{=}0.95$ (L:H\,$=\LHNinetyFive$), the mix costs
$\bar C=\CbarNinetyFive$ cycles, so utilisation 0.80 requires
$o\le\ONeededEighty$ cycles/op. That is 39\,\% below \lk{fcpq-h16-home} and
25\,\% below \lk{fc-remote}. Even $o=\OFcOld$ would give only \UtilAtOFcOld. The
ops/s gain of FC-PQ is therefore partly an accounting effect: the lock does
less work per second while completing more operations. Our first candidate
for the cost [inference] is home wakes: on 09-28, \lk{fcpq-h16} with default
placement had 769 admin cycles/op against 1\,005 for \lk{-home}. We have not
broken $o$ into drain, heap, rekey and wake costs \todo{breakdown of $o$ into
drain / heap / rekey / home-wake cycles}.

\paragraph{Usage ordering alone costs little.} At heavy ratio 1, FC-PQ's
throughput is within 0--3\,\% of FC.

\paragraph{Limitations.}
\begin{itemize}
\item Bystanders are always runnable, so steal-when-idle never fires. The
  balancing steal is our substitute, and it bounds the measurable bystander
  delay. Real bystanders are I/O-driven. With idle gaps, stealing could
  rescue trapped tasks, and it would also move a never-parking actor server.
  We have not tested this.
\item The workload is closed-loop with one await per iteration and no idle
  clients. Late joiners and clients returning from idle were not exercised,
  and no newcomer rule handles idle return.
\item Three repeats give min/max spreads, not confidence intervals.
\item The clock cap splits the data into two windows (above).
\item The workload is one data structure plus spins; there is no application
  benchmark.
\end{itemize}

% ---------------------------------------------------------------------------
\section{Related Work}\label{sec:related}

\paragraph{Scheduler-cooperative locks.} SCL and its user-space variant
U-SCL~\cite{scl} name \emph{scheduler subversion} for OS threads: lock usage
rather than scheduler share decides CPU allocation. They fix it with usage
accounting and lock-slice penalties. CFL~\cite{cfl} schedules lock occupation
by cgroup share and priority, and is NUMA-aware. Our service-fairness goal is
the same idea (fair lock \emph{time}). We apply it to a combiner in a
cooperative executor, where the combiner also decides whose thread pays.

\paragraph{Delegation locks.} FC~\cite{fc} is starvation-free and
linearizable. It lists the number of consecutive combining rounds as a knob
and does not measure the combiner's own delay. TCLocks~\cite{tclocks} bound
combining batches at 1\,024 waiters for ``long-term fairness''. They give no
fairness metric and defer charging the waiter for combining. Neither measures
a per-thread combiner share. The FC dispatcher poster~\cite{fcdispatch}
rotates the dispatcher role, which mitigates burden, but does not measure the
cost.

\paragraph{CES.} CES~\cite{ces} states that it is the first delegation-style
lock for cooperatively scheduled tasks. It evaluates throughput only, at
uniform CS cost, and notes qualitatively that the CS stays on one thread per
resource. We build on its inline/remote exchange and add what a fairness
analysis needs: chain bounds with break placement, the burden metric, and
usage-ordered service.

\paragraph{Async runtimes and locks.} The tokio mutex~\cite{tokio} is a
FIFO semaphore whose hand-off wake goes through the LIFO slot. Reading its
source, this is a bounded inline-after-poll hand-off on the unlocker's worker
[inference, consistent with the LIFO side check of \S\ref{sec:q3}]. Its fairness scope is only the waiters~\cite{tokio6049}.
\lk{async-lock}~\cite{asynclock} relies on a 0.5\,ms starvation hand-off,
which never fires when the notified waiter is never polled.

\paragraph{Preemptive request scheduling.} Shinjuku~\cite{shinjuku} preempts
every 5--15\,$\mu$s to handle heterogeneous request costs, and places contended
locks out of scope. Perséphone's DARC~\cite{persephone} reserves cores per
request type and assumes independent requests. They handle cost heterogeneity
in the request scheduler. We handle it inside the lock, where a cooperative
executor has no preemption to fall back on.

% ---------------------------------------------------------------------------
\section{Conclusion and Future Work}\label{sec:concl}

A delegation lock in a cooperative executor decides two things the executor
cannot see: which worker pays for critical sections, and which client is
served next. Left at their defaults, the first concentrates all combining on
one worker (burden Jain $1/W$) or lets one client take over, and the second
gives FIFO service proportional to cost (Jain 0.65--0.66). Both can be fixed
inside the lock with only wake-placement hints. A chain bound with home break
(\lk{ces-k64-home}) or remote wakes with a pass-end hand-off (\lk{fc-remote})
restore burden Jain 0.98--1.00 at CES throughput. Usage-ordered FC-PQ with
$H{=}16$, home wakes and a 16-pass clamp reaches service Jain 0.997 with a
17-pass wait bound. Burden-fair delegation runs 1.5--1.6$\times$
\lk{tokio::sync::Mutex} under sustained load in either tokio configuration.
Service fairness costs utilisation (0.84$\to$0.68), because the per-op lock
cost $o$ is fixed.

Future work:
\begin{itemize}
\item I/O-driven, intermittently runnable bystanders, where steal-when-idle
  exists.
\item A remote-yield actor, i.e.\ explicit server migration.
\item Clamps 9--12, to trace the predicted Jain/latency frontier.
\item Reducing $o$ below 703 cycles/op, starting with home wakes.
\item Real applications with I/O between critical sections.
\end{itemize}

{\small\bibliographystyle{plain}
\bibliography{refs}}

\end{document}
